% ============================================================
%  CONFIGURACIÓN CENTRALIZADA DEL PROYECTO (BTH / ACADÉMICO)
%  Modifique ÚNICAMENTE este archivo para personalizar todos los
%  datos de la carátula, encabezados y metadatos del documento.
% ============================================================

% ------------------------------------------------------------
% 1. DATOS INSTITUCIONALES (Normativa BTH - RM 0912/2023)
% ------------------------------------------------------------
\newcommand{\departamento}{SANTA CRUZ - BOLIVIA}
\newcommand{\distritoeducativo}{DISTRITO EDUCATIVO SAN JULIÁN}
\newcommand{\institucion}{MÓDULO TECNOLÓGICO PRODUCTIVO SAN JULIÁN}
\newcommand{\facultad}{NIVEL SECUNDARIO COMUNITARIO PRODUCTIVO} % Opcional / Universidad

% ------------------------------------------------------------
% 2. MODALIDAD Y CARRERA / ESPECIALIDAD
% ------------------------------------------------------------
% Modalidad BTH: "INNOVACIÓN TECNOLÓGICA"
\newcommand{\modalidad}{INNOVACIÓN TECNOLÓGICA}
\newcommand{\especialidad}{SISTEMAS INFORMÁTICOS}
\newcommand{\nivelformacion}{TÉCNICO MEDIO}

% ------------------------------------------------------------
% 3. DATOS DEL PROYECTO
% ------------------------------------------------------------
\newcommand{\tituloproyecto}{SISTEMA WEB DE INSCRIPCIÓN PARA EL MÓDULO TECNOLÓGICO PRODUCTIVO SAN JULIÁN BTH}
\newcommand{\subtituloproyecto}{}

% ------------------------------------------------------------
% 4. DATOS DEL / DE LOS ESTUDIANTE(S) / AUTOR(ES)
%    (El reglamento BTH permite proyectos individuales o en parejas)
% ------------------------------------------------------------
% Estudiante / Autor 1 (Obligatorio)
\newcommand{\autoruno}{Estudiante 1}
\newcommand{\emailautoruno}{}

% Estudiante / Autor 2 (Opcional: deje vacío si es individual)
\newcommand{\autordos}{Estudiante 2} % Ej: Estudiante 2 (dejar vacío si no aplica)
\newcommand{\emailautordos}{}

% Etiquetas descriptivas
\newcommand{\etiquetaautor}{Postulantes:}

% ------------------------------------------------------------
% 5. DATOS DEL TUTOR / ASESOR GUÍA
% ------------------------------------------------------------
\newcommand{\tutorproyecto}{Ing. Juan Vladimir Ramirez Flores}
\newcommand{\cargotutorproyecto}{Tutor Guía}
\newcommand{\etiquetatutor}{Tutor/a:}

% ------------------------------------------------------------
% 6. LUGAR, GESTIÓN Y FECHA
% ------------------------------------------------------------
\newcommand{\lugarproyecto}{Santa Cruz - Bolivia}
\newcommand{\gestionproyecto}{2026}
\newcommand{\fechaproyecto}{2026}

% ------------------------------------------------------------
% 7. MACROS DE COMPATIBILIDAD (Para metadatos automáticos)
% ------------------------------------------------------------
\newcommand{\autorproyecto}{\autoruno\ifx\autordos\empty\else, \autordos\fi}
\newcommand{\emailautor}{\emailautoruno}
\newcommand{\programa}{\especialidad\ (\nivelformacion)}

% ------------------------------------------------------------
% 8. TIPOGRAFÍA Y MÁRGENES DEL DOCUMENTO (docs/formato.md)
% ------------------------------------------------------------
% Tipografía del documento (docs/formato.md):
% - 'times' : Times New Roman (Oficial: 12pt, Interlineado 1.5)
% - 'arial' : Arial / Helvetica (helvet)
\newcommand{\tipografiadocumento}{times}

% Márgenes de página oficiales según docs/formato.md (adaptado para empaste/anillado):
% Izquierdo: 3.0 cm; Derecho, Superior e Inferior: 2.5 cm.
\newcommand{\margenizquierdo}{3.0cm}
\newcommand{\margenderecho}{2.5cm}
\newcommand{\margensuperior}{2.5cm}
\newcommand{\margeninferior}{2.5cm}

% ------------------------------------------------------------
% 9. ESPACIADO DE PÁRRAFOS (Estilo Microsoft Word)
% ------------------------------------------------------------
% Espaciado posterior estándar después de cada párrafo:
% - 8pt : Valor estándar predeterminado de Microsoft Word moderno (Office 2013-365)
% - 6pt : Alternativa común para formato académico compacto
% - 0pt : Formato LaTeX clásico (sin separación entre párrafos)
\newcommand{\espacioposteriorparrafo}{8pt}

% Sangría de primera línea del párrafo:
% - 0pt    : Estándar en Microsoft Word cuando hay espaciado posterior (estilo bloque)
% - 1.27cm : Estándar APA 7 tradicional (si se prefiere sangría sin espaciado posterior)
\newcommand{\sangriaprimeralinea}{0pt}

% ------------------------------------------------------------
% 9. ESPACIADOS DE PRELIMINARES Y DIAGRAMACIÓN
% ------------------------------------------------------------
% Espacio superior vertical de compatibilidad (el entorno estilodedicatoria alinea al fondo con \vfill)
\newcommand{\espaciosuperiordedicatoria}{3cm}

% ------------------------------------------------------------
% 10. CONFIGURACIÓN DE FIGURAS E IMÁGENES (NORMAS APA 7)
% ------------------------------------------------------------
% Ancho predeterminado estándar para figuras e ilustraciones:
% Se recomienda 0.8\textwidth para mantener proporciones equilibradas
% dentro de los márgenes de página carta (3.0 cm izq, 2.5 cm der).
\newcommand{\anchuraimagenpredeterminada}{0.8\textwidth}

% Espaciado vertical entre la parte inferior de la imagen y la nota al pie:
\newcommand{\espacionotafigura}{4pt}

% Estilo de presentación del rótulo (caption) para tablas y figuras:
% - 'estricto' : Norma APA 7ma Edición oficial (2 líneas: número en negrita y título abajo en cursiva).
% - 'enlinea'  : Formato adaptado compacto en 1 sola línea (ej. "Figura 1.1: Título de la figura").
\newcommand{\estilorotuloapa}{estricto}

% -------------------------------------------------------------------------
% GUÍA RÁPIDA: CÓMO INSERTAR FIGURAS BAJO NORMAS APA 7MA EDICIÓN
% -------------------------------------------------------------------------
% Según las Normas APA 7ma Edición, las figuras siguen la misma convención
% estructural y visual que las tablas:
%
% 1. RÓTULO Y TÍTULO ARRIBA DE LA IMAGEN:
%    - Se colocan con el comando estándar \caption{...}.
%    - El número de figura aparece en negrita (ej. "Figura 1.1:").
%    - El título debe ser claro, conciso y explicativo.
%
% 2. CONTENIDO VISUAL CENTRADO:
%    - La imagen debe estar centrada horizontalmente usando \centering.
%    - Las imágenes deben almacenarse en la carpeta 'imagenes/' (.png, .jpg, .pdf).
%    - Ajustar el ancho con [width=\anchuraimagenpredeterminada] o valor deseado.
%
% 3. NOTA AL PIE DE LA FIGURA (DEBAJO DE LA IMAGEN):
%    - Se coloca obligatoriamente usando la macro semántica \notafigura{...}.
%    - Imprime automáticamente el texto "Nota." en cursiva y el contenido aclaratorio.
%    - Debe indicar la fuente o autoría:
%      * De autoría propia: \notafigura{Elaboración propia.}
%      * Con base en fuentes externas: \notafigura{Adaptado de \textcite{autor2024}.}
%
% FORMAS DE INCLUSIÓN EN EL DOCUMENTO:
%
% Forma 1: Mediante el entorno flotante estándar 'figure' (Recomendado):
% \begin{figure}[htbp]
%     \centering
%     \caption{Diagrama de bloques funcional del sistema.}
%     \label{fig:diagrama_bloques}
%     \includegraphics[width=\anchuraimagenpredeterminada]{diagrama_bloques.png}
%     \notafigura{Elaboración propia con base en especificaciones técnicas.}
% \end{figure}
%
% Forma 2: Mediante la macro semántica '\incluirfigura' (Inserción rápida):
% \incluirfigura[ancho opcional]{archivo}{Título}{clave_label}{Nota al pie}
% Ejemplo:
% \incluirfigura{diagrama_proceso.png}{Flujo de proceso productivo.}{fig:proceso}{Elaboración propia.}
% -------------------------------------------------------------------------

% ------------------------------------------------------------
% 11. CONFIGURACIÓN DE PORTADA BTH (docs/portada.pdf)
% ------------------------------------------------------------
% Activar o desactivar el marco decorativo perimetral azul en la portada:
% - 'true'  : Renderiza el marco azul con esquinas ornamentadas.
% - 'false' : Portada limpia sin marco perimetral.
\newcommand{\activarmarcobth}{true}

% Recursos gráficos de la portada BTH
\newcommand{\rutamarcobth}{imagenes/marco_portada_bth.png}
\newcommand{\rutalogobth}{imagenes/logo_bth.png}

% Dimensiones de marco y logotipo institucional (docs/formato.md: 6 cm x 4.5 cm)
\newcommand{\anchomarcobth}{468.8pt}
\newcommand{\altomarcobth}{677.2pt}
\newcommand{\anchologobth}{6.0cm}
\newcommand{\alturalogobth}{4.5cm}

% Fórmula oficial de titulación académica para el grado BTH (docs/formato.md, punto 5)
\newcommand{\formulagradobth}{PROYECTO PARA OPTAR EL GRADO TÉCNICO MEDIO}
